\title{Selection-Inference: Exploiting Large Language Models for Interpretable Logical Reasoning}

\begin{abstract}
Large language models (LLMs) have been shown to be capable of impressive few-shot generalisation to new tasks. However, they still tend to perform poorly on multi-step logical reasoning problems. Here we carry out a comprehensive evaluation of LLMs on 50  tasks that probe different aspects of logical reasoning. We show that language models tend to perform fairly well at single step inference or entailment tasks, but struggle to chain together multiple reasoning steps to solve more complex problems. In light of this, we propose a Selection-Inference (SI) framework that exploits pre-trained LLMs as general processing modules, and alternates between selection and inference to generate a series of interpretable, casual reasoning steps leading to the final answer. We show that a 7B parameter LLM used within the SI framework in a 5-shot generalisation setting, with no fine-tuning, yields a performance improvement of over 100\% compared to an equivalent vanilla baseline on a suite of 10  logical reasoning tasks. The same model in the same setting even outperforms a significantly larger 280B parameter baseline on the same suite of tasks. Moreover, answers produced by the SI framework are accompanied by a \emph{causal} natural-language-based reasoning trace, which has important implications for the safety and trustworthiness of the system.
\end{abstract}

\section{Introduction}
Large language models (LLMs) are powerful few-shot learners \citep{bommasani2021opportunities, brown2020language, lu2021pretrained}. However, one area where they tend to perform poorly is logical reasoning \citep{rae2021scaling}. Yet the ability to perform multi-step, logically valid reasoning is fundamental for the discovery of new knowledge and explainability. It underpins many advancements that have been made in science, medicine, maths and philosophy. It is also one of the most valued strengths of classical, symbolic AI over contemporary deep learning methods \citep{marcus2019rebooting, marcus2020next,bengio2021deep}, prompting the recent increase in the use of neurosymbolic approaches to bridge this gap \citep{garnelo2019reconciling,garcez2020neurosymbolic}. Here we propose a Selection-Inference (SI) framework that takes inspiration from the neurosymbolic literature to improve the ability of LLMs to do logically valid reasoning.

There are many flavours of neurosymbolic models \citep{garcez2020neurosymbolic}. Those from which we draw inspiration tend to have a modular structure, where each module is specialised for one type of operation \citep{mao2019neuro, andreas2016neural}. For example, such modules may be neural networks or hand-crafted functions designed to attend to a single object, or to compare the location or size of two inputs \citep{andreas2016neural, yi2018neural}. Neurosymbolic models can produce an answer to a complex query by chaining these operations together, passing inputs from one module to another. This has the benefit of producing an interpretable trace of intermediate computations, in contrast to the ``black-box'' computations common to end-to-end deep learning approaches. Importantly, the modularity of neurosymbolic methods allows them to generalise to significantly harder problems that require long chains of reasoning \citep{hudson2019learning}. However, the hand-crafted and specialised nature of the modules often makes the resulting systems brittle, hard to optimise, and difficult to extend to new domains \citep{yi2018neural}.

Our SI framework, drawing on best practice from neurosymbolic approaches, decomposes logical reasoning into two modular stages: 1) \emph{selection}, which involves choosing a subset of relevant information sufficient to make a single step of inference, and 2) \emph{inference}, which only sees the limited information provided by the selection module, and uses it to infer a new intermediate piece of evidence on the way to the final answer (see Fig.~\ref{fig:SI_framework}c). We implement both stages using pre-trained LLMs which, thanks to their powerful few-shot generalisation capabilities, serve as more general alternatives to the hand-crafted, specialised modules typically used in neurosymbolic approaches. In the SI framework, multiple steps of selection and inference are chained together to produce a sequence of reasoning steps. As well as underpinning better performance on reasoning problems, this yields an interpretable trace that justifies the final answer.

Furthermore, the reasoning trace produced by our system is \emph{causal}, in the sense that each step follows from, and depends on, the previous step. Each inference step is made in isolation, based solely on the limited information provided by the Selection module, without direct access to the question or to previous steps of reasoning. This contrasts with the more common approach of obtaining \emph{post-hoc rationalisation}, where the answer produced by the model has no direct dependence on the explanation, since the explanation is produced either in parallel to the answer or after the fact \citep{saha2020prover, lampinen2021tell, cobbe2021training}. A notable example that sits in the grey area between post-hoc rationalisation approaches and the more causal explanation approaches is Chain-Of-Thought (COT) \citep{wei2022chain} (see Fig.~\ref{fig:SI_framework}b). In this approach LLMs are encouraged to produce a reasoning trace before the answer. However the dependence of the answer on the reasoning is not explicitly encouraged to be causal (as defined above). Indeed, the authors show that while the COT explanations help boost the final answer accuracy, the reasoning traces produced by the model are often wrong even when the final answer is correct (see the Supplementary Materials of \citep{wei2022chain} for examples). Developing a system that can demonstrate how it reaches its answers using a \emph{causal} reasoning trace has important benefits in terms of safety, explainability, interpretability, debugging, and trust. In this paper we make the following contributions:

\begin{enumerate}
    \item We provide a comprehensive evaluation of LLMs on a set of 50  tasks probing different aspects of logical reasoning, and show that LLMs are good at simpler single step logical inference in 5-shot generalisation settings, but struggle with harder problems (Sec.~\ref{sec:how_well_do_llms_reason})
    \item We introduce the Selection-Inference (SI) framework, a modular, iterative approach to solving reasoning problems (Sec.~\ref{sec:selection_inference}).
    \item We demonstrate the utility of the SI framework by evaluating a 7B parameter LLM from the Gopher family \citep{rae2021scaling} on 10  logical reasoning tasks, showing overall that it almost triples the performance of the same model used naively and almost doubles the performance of the same model used in the COT framework. Moreover, it often outperforms a 40x larger 280B LLM baseline used both naively and in the COT framework. 
    \item We illustrate further benefits of the SI framework in terms of the causal and interpretable reasoning traces produced (Sec.~\ref{sec:si_reason_traces}). These traces can help humans understand how the model reached its final answer, which is useful for debugging and opens the system's decisions to human critique.
\end{enumerate}

\section{Related Work} Our Selection-Inference framework sits at the intersection of classical, symbolic AI and deep learning. A typical symbolic AI system might consist of a knowledge base, which is typically hand-curated by experts, and an inference engine that allows the system to perform logic-based reasoning over its knowledge base. For example, such a system could apply step-by-step reasoning to answer a complex question such as ``What are the apothecary's incentives and disincentives for selling poison to Romeo in Romeo and Juliet?'' \citep{lenat_2019} - something that even the best contemporary deep learning based systems struggle to do.

One of the primary benefits of symbolic AI systems over deep learning models is their interpretability; we can look at the reasoning steps such a system has taken to see how the final conclusion was reached. However, unlike deep learning approaches, symbolic AI systems require knowledge to be hand-crafted and are in general hard to scale. Although some approaches have attempted to bridge the gap between deep learning and symbolic AI by converting problems into formal logic \citep{nye2021improving} and using existing solvers to help produce an answer, this process can be brittle and tends not to scale well. Another attempt to bridge the gap comes from the neurosymbolic perspective \citep{mao2019neuro, yi2018neural, gupta2019nmn, hudson2019learning}. These models combine the best parts of deep learning -- learning knowledge from data -- and symbolic AI -- producing an interpretable reasoning trace. However, they are typically quite brittle due to the hand-crafted \citep{gupta2019nmn}, specialised nature \citep{mao2019neuro} of the modules and optimisation difficulties.

On the deep learning side, recent work has attempted to adapt large pre-trained language models, LLMs, to the task of logical reasoning. At a high level these can be split into three groups: 1) approaches that try to fine-tune LLMs to produce the final answer directly, keeping reasoning implicit \citep{betz2020critical, clark2020transformers} (e.g. Fig.~\ref{fig:SI_framework}a); 2) approaches that encourage LLMs to produce reasoning explicitly, but all reasoning steps are produced in one generative step \citep{nye2021show, zelikman2022star,jhamtani2020learning,Dalvi2021ExplainingAW,wei2022chain,cobbe2021training} (e.g. Fig.~\ref{fig:SI_framework}B); and 3) approaches that use LLMs to produce each reasoning step one at a time \citep{tafjord2020proofwriter}. The latter is where our Selection-Inference framework sits (Fig.~\ref{fig:SI_framework}C). In general it was found that the approaches that incorporate explicit reasoning work better than those that only try to predict the final answer. However, although explicit reasoning helps improve the accuracy of the models, encouraging the models to produce multiple steps of reasoning in a single generative pass is not enough to make the models use reasoning in a causal manner. The authors found that the generated reasoning traces often contain unrelated or incorrect steps while still resulting in the correct answer (see examples in the appendices of \citep{zelikman2022star, wei2022chain}). Encouraging LLMs to generate each reasoning step one at a time \citep{tafjord2020proofwriter} is currently the most promising direction for achieving causal reasoning, and it is the approach we take in our paper. While the model proposed by \cite{tafjord2020proofwriter} is very impressive, it only works for ``Prove this statement to be True/False'' style questions, since it relies on enumerating all possible implications and checking whether the question statement or its negation are present in the inferred facts, which is also computationally expensive. 

\section{How Well Do Large Language Models Reason?}\label{sec:how_well_do_llms_reason}

Past work has shown that LLMs are poor at logical reasoning \citep{rae2021scaling}, however the evaluation was done on a relatively small set of tasks, and was not systematic. In particular, here we are interested in 1) how LLMs perform on simple entailment tasks compared to multi-step reasoning problems and 2) how scaling laws -- suggested by \cite{rae2021scaling} -- apply to logical reasoning. To this end, we evaluated LLMs from the Gopher family in a 5-shot\footnote{We chose 5-shot setting because it was used in \citep{rae2021scaling}, and because \citep{min2022rethinking} have demonstrated that additional shots beyond 4 result in limited increase in multi-choice accuracy} setting on a larger set of 50  tasks that touch on different aspects of logical reasoning and vary in terms of the number of reasoning steps required, presence or absence of negation, whether the relevant context information was provided, and whether the model is required to evaluate the accuracy of multiple choices or generate the answer among others. The additional tasks were collected from six sources: bAbI \citep{weston2015towards}, BigBench \citep{BigBench}, AAC \citep{betz2020critical}, Jeopardy \citep{tunguz_2019}, ProofWriter \citep{tafjord2020proofwriter} and 2WikiMultiHop \citep{welbl2018constructing} (see Fig.~\ref{fig:logic_baselines} in Supplementary Information for raw results).

Our analysis found that LLMs are good at some aspects of logical reasoning (Fig.~\ref{fig:baseline_multistep_reasoning}). In particular, they appear to be good at simple entailment and implication tasks (e.g. see AAC tasks and Entailed Polarity in Fig.~\ref{fig:logic_baselines}). This appears to hold when negation is present (AAC Split Extended tasks), and both in generative (AAC Split) and multiple-choice scoring settings (AAC Split Extended tasks, Entailed Polarity). However, the performance of vanilla language models tends to decrease when they get presented with irrelevant facts alongside the ones relevant for reasoning (e.g. see 2WikiMultiHop With Context tasks, bAbI tasks 2-3 or ProofWriter tasks), when they have to infer the relevant facts from memory (e.g. 2WikiMultiHop or StrategyQA tasks), and as the questions start to require more steps of reasoning (e.g. see the performance drop between bAbI tasks 1-3 or ProofWriter Tasks).

In line with Rae et al.'s findings, our results confirmed that LLMs of larger sizes do perform better than the smaller models. However, we found that even the largest 280B model performed only 13.6\% above chance level on average across the 38 available multi-choice logic tasks  (see Figs.~\ref{fig:logic_baselines}-\ref{fig:logic_baselines_multichoice_less_random} and Sec.~\ref{si_sec:baseline_results} in Supplementary Information for more details). Furthermore, we found that logical reasoning tasks were qualitatively different from other natural language tasks. The scaling law for logic-based tasks within the Gopher family models was significantly worse than for other language tasks measured here as the average performance on the subset of BigBench tasks from \citep{rae2021scaling} with the logic tasks used in this paper removed (see Fig.~\ref{fig:scaling_laws}). 

\section{The Selection-Inference (SI) Framework} \label{sec:selection_inference}

We are interested in solving logical reasoning problems expressed in natural language. In this work we assume that each question is accompanied by context information (see Figs.~\ref{fig:SI_framework} and \ref{fig:si_example_babi}), which contains all the information necessary to solve the problem, as well as potentially irrelevant distractors. In the future this assumption can be relaxed, for example by extracting the necessary information through search \citep{lazaridou2022internet, menick2022teaching}. We also assume that all questions are well posed and definitively answerable given the context. 

Logical reasoning problems require using existing information to infer new relevant knowledge necessary to answer the question. This can be done through deduction, induction or abduction, although the datasets we use here contain mostly deductive and a small number of inductive problems\footnote{See Fig.~\ref{fig:si_qualitative_results} for an example of deduction and induction problems used in this paper.}. Some problems require multiple steps of inference, where later steps use the knowledge inferred in the earlier steps. Hence, we use an iterative framework where at each step the SI uses information in the existing context, $\mathcal{C}_t$, to infer a new fact, $f_t$, which is appended back to the context to create new context, $\mathcal{C}_{t+1} = \mathcal{C}_t \cup f_t$. This process can then iterate until the solution to the question is found. In the current implementation of the SI framework, we repeat the process for a fixed number of steps and take the final inference to be the answer. Addressing the issue of halting is left for future work.

Inspired by neurosymbolic methods, we additionally split each step of reasoning into further two components: 1) \emph{Selection}, which selects a subset of the information present in the context, $s_t$, given the context and the question, $\mathcal{C}^t \cup q$. This selection, $s^t$ is fed to the next step, 2) \emph{inference}, which produces the new fact, $f_t$, based on the information passed to it by the selection step. Examples of selection and inference are shown in Fig.~\ref{fig:si_example_babi}. This splitting of each step of reasoning into selection and inference is the main contribution of our paper, and is important for several reasons. First, and most importantly, it makes the resulting reasoning \emph{causal}, since both steps have limited capabilities by design, and are interdependent. The selection step is constrained to only use the information available in the context, an the inference step only sees the subset of facts provided by the selection without access to the question. Hence, the model is unlikely to make up information to answer the question, and it cannot ignore reasoning when producing the final answer. The other benefit of this approach is that each step of reasoning is broken down into even smaller sub-tasks, which are easier for LLMs to adapt to, and which helps make the reasoning more generalisable to harder problems. 

In this paper we use pre-trained, frozen language models from the Gopher family \citep{rae2021scaling} in a 5-shot generalisation setting using prompt engineering to implement the Selection and Inference modules. We settled on prompt engineering to evaluate the base utility of the SI framework, however it can also be used in the fine-tuning setting which we explore briefly in Sec.~\ref{sec:fine_tuning}. We next describe the Selection and Inference modules in more detail.

\subsection{Selection Module}\label{sec:selection_module}

We use prompt engineering to encourage the model to output the correct selection, $s^t$. The n-shot prompt is a string of the the following form:

\begin{small}\begin{lstlisting}
"""
# n-shot prompt # First example. <context 1> <question 1> # Example selection <fact>. We know that <fact>[ and <fact>]*. Therefore, ... # Problem to solve. <context> <question>

"""
\end{lstlisting}\end{small}

where \verb|<fact>|s are copied directly from the \verb|context|, and \verb|[ and <fact>]*| means that the module is allowed to select more than one fact for each step of inference, where the total number of facts is a hyper-parameter. 

The simplest option to implement the selection is to feed this prompt directly to a pre-trained LLM and take the output generated by the language model. However, this unconstrained approach may allow the model to make up facts, thus removing the causal aspect of the reasoning trace. Indeed during experimentation this is what we often found. So instead we use the pre-trained LLM to score each of the facts in the context, and select the one with the highest log-likelihood. We can then repeat this process by appending each new fact to the end of the previous prompt until the full selection is constructed. Note that for now we \texttt{halt} after a fixed number of steps. See Algorithm \ref{alg:scoring_selection} for more details.

\subsection{Inference module} \label{sec:inference}

The n-shot prompt for the Inference module has the following form (shown below, also see Fig.~\ref{fig:SI_framework}):

\begin{small}\begin{lstlisting}
    """
    #n-shot inference prompt # First example. <fact>. We know that <fact>[ and <fact>]*. Therefore, <new fact>. ...

    # Problem to solve. <output of the Selection module>. Therefore, """
\end{lstlisting}\end{small}

The n-shot prompt and the output of the Selection module, are fed to the pre-trained LLM serving as the Inference module. The first generated sentence (extracted from the generated text as per BigBench evaluation \citep{BigBench} pipeline, see Supplementary Materials for details) is taken to be the newly inferred fact. This fact is added to the context, which concludes one reasoning step of the SI framework. For now, we halt after a fixed number of steps. See Algorithm \ref{alg:selection_inference} for more details.

\begin{small}\begin{algorithm}
\caption{Selection-Inference}\label{alg:selection_inference}
\begin{algorithmic}
\Require \texttt{An n-shot selection prompt,} $p_{select}$.
\Require \texttt{An n-shot inference prompt,} $p_{infer}$.
\Require \texttt{Initial Context,} $\mathcal{C}^0$, \texttt{made up of facts (and rules).}
\Require \texttt{The question,} $q$.
\Require \texttt{Language model, LLM.}
\Require \texttt{A halting function, halt, determines if the answer has been reached.}
\State $t=0$ \Comment{Start at step $0$}.
\While{\texttt{not halt()}}
    \State $s^t \gets$ \texttt{Selection\_Module}$(p_{select}, C^t, q, \texttt{LLM})$ \Comment{Do selection.}
    \State $i^t \gets$ \texttt{Inference\_Module}$(p_{infer}, s^t)$ \Comment{Do inference.}
    \State $C^{t+1} \gets C^t \cup i^t$ \Comment{Add the newly inferred fact to the context.}
    \State $t \gets t+1$ \Comment{Move onto the next step of reasoning}

\EndWhile \\
\Return{$s^t$}
\end{algorithmic}
\end{algorithm}
\end{small}

\section{Experiments and Results}
\label{sec:experiments_results}
We evaluate our SI framework on a subset of 10 /50  logical reasoning tasks introduced in Sec.~\ref{sec:how_well_do_llms_reason}. These tasks were chosen based on whether they include context information necessary to answer the question, whether the questions have a definitive answer, and to ensure that they cover different kinds of reasoning abilities. The tasks include bAbI \citep{weston2015towards} Tasks 1-3, which require the model to use 1-3 supporting time-ordered facts respectively to answer a question, and Tasks 15-16, which test deductive and inductive reasoning respectively. We also evaluate our model on the ProofWriter OWA datasets \citep{tafjord2020proofwriter} of depth 0, 1, 2, 3 and 5 (there is no depth 4 task). These are language-based logical reasoning problems, where the depth is the number of reasoning steps required to answer the question.

To baseline the performance of the SI framework, we consider a 7B (same size as the LLM used in the SI framework) and a 40x larger 280B parameter LLM evaluated in a 5-shot setting. There are two types of evaluation for these vanilla baselines that we consider: \textit{multi-choice} and \textit{generative} evaluation. In \textit{generative} evaluation, we measure the exact string match (first sentence in lower case and ignoring any non-alphabetic characters) between the output generated by the LLM and the ground truth answer. This is appropriate, since most of the tasks that we consider require either a single word answer, or the dataset is such that the answers are highly structured. In \textit{multi-choice} evaluation the LLM is used to score each of the answer choices, as in Li et al. \citep{li2021systematic}. In general LLMs perform significantly better in a \textit{multi-choice} vs \textit{generative} evaluation setting, since the chance level in the multi-choice setting is significantly higher.

We also consider a chain-of-thoughts (COT) \citep{wei2022chain} inspired baseline, where the k-shot prompts to the 7B and 280B models include reasoning traces for the same examples that we use to prompt the SI framework (although with selection and inference combined, see Supplementary Information for example prompts). This tests whether providing the reasoning examples alone is sufficient to improve performance, or whether the further breakdown into Selection and Inference sub-steps improves accuracy. Note that among all of the approaches outlined only the SI framework is explicitly set up to generate \emph{causal} reasoning traces. 

Fig.~\ref{fig:summary_results} demonstrates that overall when evaluated generatively, the 7B parameter LLM within the SI framework performs better (58.75\%) than the same model evaluated naively (2.94\%) or in the COT framework (41.32\%) (all $p<0.01$, see Supplementary Information for details of the calculations). Not only that, the 7B parameter LLM within the SI framework also outperforms on average the 40x larger 280B parameter LLM in both vanilla (31.19\%) and COT framework (44.03\%) (all $p<0.01$). When evaluated in the easier multi-choice setting, we find that surprisingly\footnote{This \emph{could} suggest that the 280B LLM has stronger priors, than the 7B LLM, which it favours over logical reasoning. For example, favouring \texttt{sheep are afraid of wolves} despite a context to the contrary \citep{min2022rethinking}. However this requires further investigation.} the vanilla 7B parameter LLM outperforms the 280B parameter LLM (57.31\% vs 51.45\%), while still performing significantly worse than the 7B SI model ($p=0.012$). Note that the latter is evaluated in the harder generative setting. Per task breakdown shown in Fig.~\ref{fig:quant_results} demonstrates that the SI framework solves the bAbI 15 deduction task, the only model to achieve 100\% accuracy (significant difference from the other models, $p<0.01$). Furthermore, it does so having seen only five examples in the prompt. The 7B SI model also significantly outperforms all other models on ProofWriter Depth 0 ($p<0.01$), ProofWriter Depth 1 ($p=0.034$).

As well as improving upon most baselines quantitatively, the SI framework also has additional qualitative benefits: 1) it produces a causal, human interpretable reasoning trace that shows how the model reached its answer and 2) it is able to recover from errors. We will now discuss each of these in turn. 

Since the Selection module is only allowed to pick facts from the context and is separate from the Inference module, and since the latter does not have access to the question, the model has to use what is selected and cannot bypass the selection to compute its answer, thus creating a causal reasoning trace. Since the reasoning trace is in natural language and is causal, it can be audited and inspected by humans, which has significant implications for safety and interpretability. 

Example reasoning traces produced by the SI model are shown in Fig.~\ref{fig:si_example_babi}. In the bAbI 16 example shown on the right the model is solving an inference problem, which requires the model to infer the colour of an animal given facts about the colours of other animals. In this example, the model is asked \texttt{"What colour is greg"}, and told that \texttt{"greg is a lion"}. This means first the model needs to use induction to infer a rule about the colour of lions. On the first step, we see that the model induces a rule, \texttt{"lions are white"}, based on the fact that \texttt{"brian is a lion"} and \texttt{"brian is white"}; we can see exactly what data underlies the model's decision to form a new rule. On the second step, we see that the model applies this newly inferred rule to the fact that \texttt{"greg is a lion"} to reach the final conclusion that \texttt{"greg is white"} using deduction. Note that the ability of the SI framework to produce inductions relies on its ability to deal with uncertainty and understand what is ``reasonable'' - something that LLMs are naturally capable of, while also being a known struggle point for symbolic AI. 

Since the reasoning traces are output in natural language, they are easy for humans to interpret and potentially intervene. Consider a scenario where there may be both a white lion and a green lion mentioned in the context, in which case we could see which information the model used to make its final decision and decide whether we want to trust it (example in Fig.~\ref{fig:recover_from_error}). We could also imagine examples where the model puts together two unrelated facts to come up with an incorrect inference, and this could also be easily be spotted by a human and rectified by replacing the wrongly inferred fact with a correct one, and re-running the consequent reasoning steps. 

Aside from inspecting reasoning traces and using them to debug when something goes wrong, the additive nature of our model - it accumulates new knowledge with each reasoning step, means that it also has the ability to recover from errors. Fig.~\ref{fig:recover_from_error} demonstrates such an example. In the first step the model inferred that \texttt{"swans are often gray"}, using the facts that \texttt{"julius is a swan"} and \texttt{"julius is gray"}. While this is correct, this new information is not useful for answering the question, which asks about lions. However, it is still possible for the model to make the more useful inference that \texttt{"lions are often white"} in a later step and recover from its original misstep.

\section{Fine-tuning Language Models for Selection and Inference} \label{sec:fine_tuning}
In Sec.~\ref{sec:experiments_results} we have demonstrated significant improvements in logical reasoning accuracy when using prompt-engineering to specialise LLMs for Selection and Inference in the SI framework. Prompt-engineering has the additional benefit of not requiring large amounts of step-by-step reasoning data, which may be hard to obtain. In this section we investigate whether the accuracy of the SI framework can be further improved by fine-tuning the LLMs for Selection and Inference. We use the ProofWriter dataset for which ground truth reasoning traces exist. 

The Selection LLM is fine-tuned to select a subset of sentences (including facts and rules) from the context by generating a string of the form \emph{"sent 2. We know that sent 4 [and sent 7]*."} given the context and the question. We ask the Selection LLM to generate sentence labels (e.g. \emph{"sent 2"} or \emph{"sent 4"}) instead of the sentences themselves, because this prevents the Selection LLM from cheating by making up facts to answer the question quicker. This preserves the dependency of the selection and therefore subsequent reasoning steps on the context. The Inference LLM is fine-tuned to compute an entailment given the selection. Both models are fine-tuned on single steps of reasoning only. Example training data are shown in Fig.~\ref{fig:fine_tune_train_data}. 

The Inference LLM converged very quickly to >99\% test accuracy after only 300 fine-tuning steps with a batch size of 16, which is to be expected as we found that pre-trained LLMs are good at single step entailment out of the box as shown in Fig.~\ref{fig:baseline_multistep_reasoning}. Examples of inferences made by the model are shown in Fig.~\ref{fig:fine_tune_inf}. The Selection LLM was trained for $4\times10^4$ steps (with batch size 16 for 50 hours on a TPU) with the exact string match accuracy reported in Fig.~\ref{fig:fine_tune_select}. Although we notice that the model is much better at predicting selections for problems that require fewer steps of inference than those that require more, ultimately the model still achieves high ($>80\%$) accuracy across most of the reasoning depths. Predicting early selections for deeper reasoning problems is hard, because it requires planning. It is an important problem to address in future work.

Fig.~\ref{fig:quant_results} shows that fine-tuning LLMs on single steps of reasoning within the SI framework leads to significant improvements in final reasoning accuracy (78.95\%) over the prompt-engineered version of the SI framework (57.93\%) and other prompt-engineered baselines (vanilla/COT generative 7B: 0.34/15.73\%, 280B: 31.58/21.12\%). We also found that the fine-tuned 7B LLM used within the SI framework produces significantly more accurate reasoning traces compared to the same LLM fine-tuned to predict all reasoning steps in one go (Fig.~\ref{fig:fine_tune_vs_baseline}). We quantified this using the Jaccard Similarity, $\texttt{Jaccard Similarity} = (M \cap GT) / (M \cup  GT)$, between the proof steps predicted by each model, $M$, and the ground-truth reasoning steps, $GT$, as shown in, calculated using exact string match over alphanumeric characters.

Qualitatively we observed that while the baseline model is good at predicting most of the reasoning steps, they often appear in the wrong order, there are additional reasoning steps that are not on the minimal reasoning path, and some steps get repeated a number of times. 

\section{Conclusion}\label{sec:conclusion}
We have presented the Selection-Inference framework for improving the ability of pre-trained language models to solve logical reasoning problems expressed in natural language. Our approach borrows from the best practices of neurosymbolic approaches to break down logical reasoning into a modular recursive pipeline that not only significantly improves the reasoning accuracy, but also produces a \emph{causal} interpretable reasoning trace. We have demonstrated that prompt-engineered LLMs used in the SI framework significantly outperform both the vanilla and COT baselines evaluated in equivalent settings, and even 40x larger baselines. The performance of the SI framework can be further improved through fine-tuning if step-by-step reasoning data is available.

A model capable of casual, interpretable and logically valid multi-step reasoning has potential applications in law, medicine, science, maths, economics, and other areas where trustworthy and verifiable logical inference is important. At the same time we recognise that special care will need to be taken to evaluate such a model before deployment in any of these settings. Further work is also needed, for example, to improve the Selection module (e.g. by allowing the model search over and evaluate different reasoning traces); to address the halting issue (both in terms of when to stop the selection and when to stop the overall reasoning); to incorporate verifiers that would help avoid false inferences being added to the context; to enable the system to source its own relevant context rather than relying on it being provided in the dataset; and to extend the ability of the system to deal with ambiguous or unanswerable questions.

\end{document}
